\documentclass[11pt,a4paper]{article}

% ── Codificación y fuentes ────────────────────────────────────────────────────
\usepackage[utf8]{inputenc}
\usepackage[T1]{fontenc}
\usepackage{lmodern}
\usepackage[spanish]{babel}

% ── Geometría y espaciado ─────────────────────────────────────────────────────
\usepackage[top=2.5cm, bottom=2.5cm, left=2.8cm, right=2.8cm]{geometry}
\usepackage{setspace}
\setstretch{1.15}
\usepackage{parskip}

% ── Tablas ────────────────────────────────────────────────────────────────────
\usepackage{booktabs}
\usepackage{tabularx}
\usepackage{array}
\usepackage{enumitem}

% ── Colores y cajas ───────────────────────────────────────────────────────────
\usepackage[dvipsnames]{xcolor}
\usepackage{tcolorbox}
\tcbuselibrary{skins, breakable}

% ── Cabeceras y pies de página ────────────────────────────────────────────────
\usepackage{fancyhdr}
\usepackage{lastpage}
\pagestyle{fancy}
\fancyhf{}
\fancyhead[L]{\small\textbf{Informe de Análisis} --- Imágenes}
\fancyhead[R]{\small 2026-06-27}
\fancyfoot[C]{\small Página \thepage\ de \pageref{LastPage}}
\renewcommand{\headrulewidth}{0.4pt}

% ── Hiperenlaces y URLs ──────────────────────────────────────────────────────
\usepackage[colorlinks=true, linkcolor=NavyBlue, urlcolor=NavyBlue]{hyperref}
\usepackage{url}

% ── Permitir saltos de línea flexibles (evita desbordes de margen) ────────────
\sloppy
\emergencystretch 3em

% ── Colores personalizados ────────────────────────────────────────────────────
\definecolor{azulTitulo}{HTML}{1B2A6B}
\definecolor{grisFondo}{HTML}{F5F5F5}

% ── Estilos de cajas ─────────────────────────────────────────────────────────
\tcbset{
  cajaTitulo/.style={
    enhanced, breakable,
    colback=azulTitulo!8, colframe=azulTitulo,
    fonttitle=\bfseries\large, coltitle=white,
    attach boxed title to top left={yshift=-2mm, xshift=4mm},
    boxed title style={colback=azulTitulo},
    top=6pt, bottom=6pt, left=8pt, right=8pt,
  },
  cajaContenido/.style={
    enhanced, breakable,
    colback=grisFondo, colframe=gray!50,
    fonttitle=\bfseries, coltitle=black,
    top=4pt, bottom=4pt, left=8pt, right=8pt,
  },
}

% ─────────────────────────────────────────────────────────────────────────────
\begin{document}

% ══ PORTADA ══════════════════════════════════════════════════════════════════
\begin{center}
  {\Large\bfseries\color{azulTitulo} ELECTRÓNICA --- Análisis de Imágenes}\\[4pt]
  {\large Informe de Análisis}\\[12pt]
  \rule{\linewidth}{1.2pt}\\[6pt]
  \rule{\linewidth}{0.6pt}
\end{center}

% ══ FICHA TÉCNICA ════════════════════════════════════════════════════════════
\vspace{4pt}
\begin{tcolorbox}[cajaTitulo, title=Datos del análisis]
\begin{tabular}{@{}llll@{}}
  \textbf{Fecha:}       & 2026-06-27  &
  \textbf{Modelo usado:} & gemini-3.5-flash \\[4pt]
  \textbf{Archivos:}    & \multicolumn{3}{l}{\texttt{imagen\_1.jpeg}, \texttt{imagen\_2.jpeg}} \\
\end{tabular}
\end{tcolorbox}

% ══ DESCRIPCIÓN GENERAL ══════════════════════════════════════════════════════
\section*{Descripción general}
Dos capturas de pantalla de la interfaz de Seguridad de Windows que muestran alertas de seguridad de nivel de gravedad 'Alta'. El sistema ha detectado y bloqueado herramientas de hacking basadas en Python (Impacket) dentro de un entorno virtual de desarrollo dedicado a la ciberseguridad.

% ══ ELEMENTOS DETECTADOS ═════════════════════════════════════════════════════
\section*{Elementos detectados}
\begin{itemize}[leftmargin=1.5em, itemsep=2pt]
  \item \textbf{Alerta de Amenaza Detectada (Imagen 1)}: Muestra una amenaza clasificada como
    \texttt{HackTool:Python/Impacket.Q} en estado ``En cuarentena''. El archivo afectado es
    \texttt{GetNPUsers.py} ubicado en la ruta del entorno virtual de ciberseguridad.
  \item \textbf{Alerta de Amenaza Activa (Imagen 2)}: Muestra una amenaza clasificada como
    \texttt{HackTool:Python/Impacket!AMTB} en estado ``Activo''. Indica que no se corrigieron
    las amenazas activas y se están ejecutando. Los archivos afectados son
    \texttt{spnego.py} e \texttt{\_\_init\_\_.py} dentro de la librería \texttt{impacket}.
  \item \textbf{Ruta de archivos afectados}:

    \path{C:\Users\E550\Desktop\CYBERSEGURIDAD\venv\...}
\end{itemize}

% ══ TEXTO EXTRAÍDO ═══════════════════════════════════════════════════════════
\section*{Texto extraído}
\begin{tcolorbox}[cajaContenido, title=Transcripción de texto visible en las imágenes]
{\small
Todos los elementos recientes --- Filtros

\textbf{Amenaza encontrada: se requiere acción.}\hfill 27/6/2026 10:32 a.\,m. --- Alta

\textbf{Amenaza puesta en cuarentena}\hfill 27/6/2026 10:31 a.\,m. --- Alta

\medskip
Detectado: \texttt{HackTool:Python/Impacket.Q}\hfill Estado: En cuarentena\\
Los archivos en cuarentena están en un área restringida donde no pueden dañar el dispositivo.
Se quitarán automáticamente.\\
Fecha: 27/6/2026 10:32 a.\,m.\\
Detalles: Este programa tiene un comportamiento potencialmente no deseado.\\
Elementos afectados:\\
\path{file: C:\Users\E550\Desktop\CYBERSEGURIDAD\venv\Scripts\GetNPUsers.py}

\medskip
Detectado: \texttt{HackTool:Python/Impacket!AMTB}\hfill Estado: Activo\\
No se corrigieron amenazas activas y se están ejecutando en el dispositivo.\\
Fecha: 27/6/2026 10:32 a.\,m.\\
Detalles: Este programa tiene un comportamiento potencialmente no deseado.\\
Elementos afectados:\\
\path{file: C:\Users\E550\Desktop\CYBERSEGURIDAD\venv\Lib\site-packages\impacket\spnego.py}\\
\path{file: C:\Users\E550\Desktop\CYBERSEGURIDAD\venv\Lib\site-packages\impacket\__init__.py}
}
\end{tcolorbox}

% ══ OBSERVACIONES ════════════════════════════════════════════════════════════
\section*{Observaciones}
Las detecciones corresponden a la suite 'Impacket', una colección de clases Python para trabajar con protocolos de red, ampliamente utilizada en pruebas de penetración (pentesting) y ciberseguridad. Dado que la ruta del archivo contiene la carpeta 'CYBERSEGURIDAD', es muy probable que se trate de un entorno de aprendizaje o desarrollo intencional, por lo que estas alertas representan falsos positivos en el contexto del usuario, aunque son firmas legítimas de herramientas de hacking para el antivirus.

% ══ SUGERENCIAS ═════════════════════════════════════════════════════════════
\section*{Sugerencias}
Si el uso de estas herramientas es legítimo y controlado (por ejemplo, para prácticas de ciberseguridad o desarrollo):
\begin{enumerate}[leftmargin=2em, itemsep=2pt]
  \item Acceda a la opción ``Acciones'' en cada alerta de Seguridad de Windows.
  \item Seleccione ``Permitir en el dispositivo'' o restaure los archivos de la cuarentena.
  \item Para evitar futuras interrupciones, agregue una exclusión en Seguridad de Windows
    para la carpeta contenedora:\\[2pt]
    \path{C:\Users\E550\Desktop\CYBERSEGURIDAD}
\end{enumerate}
Si no reconoce estas herramientas, proceda a eliminarlas inmediatamente a través del botón
``Acciones'' $\rightarrow$ ``Quitar''.

% ══ PIE DE INFORME ════════════════════════════════════════════════════════════
\vspace{12pt}
\begin{center}
  \footnotesize\color{gray}
  Informe generado automáticamente por el sistema de análisis con IA (gemini-3.5-flash) y soporte LaTeX. \\
  Generado el 2026-06-27 \textbullet\ ELECTRÓNICA --- Sistema de Análisis de Imágenes.
\end{center}

\end{document}
